\section{\emph{Summorum Pontificum} and the Normalization of Demand}

Benedict XVI's \emph{Summorum Pontificum}, issued 7 July 2007 and effective 14 September, changed the general legal environment in which the FSSP worked. It called the 1962 Missal an extraordinary expression of the Roman Rite and replaced the earlier indult conditions. A Latin priest could use it in Mass without a congregation without permission; stable groups acquired a defined route through pastors, bishops, and the Pontifical Commission; communities of institutes and societies could provide for community use under their superiors and proper law; and local ordinaries could erect personal parishes.

For the FSSP, the act did not supply an identity that had been missing. The society already possessed particular founding decrees and definitive constitutions. The change was contextual: priests and faithful outside \emph{Ecclesia Dei} institutes received broader access, dioceses encountered more requests, and personal-parish structures gained an explicit place in the general legislation. The FSSP's trained clergy and common life made it a ready partner for bishops responding to that demand.

The period saw continued geographic expansion, seminary enrollment, ordinations, and personal parishes. Growth did not mean that every request yielded an apostolate. A bishop still decided whether to receive a house or entrust ministry, and a finite number of priests had to cover existing communities. Nor did \emph{Summorum Pontificum} make the reformed Missal invalid or marginal in law; article 1 called it the ordinary expression of the Latin Church's \emph{lex orandi} under the categories then used.

\subsection{A change in social function}

In 1988 the new Fraternity served as one solution for clergy and faithful separating from Lefebvre's act. By the 2007 regime, later cohorts included Catholics too young to have personal memory of Écône or the immediate postconciliar transition. The checked public set contains no member survey showing how candidates encountered the older books or why they entered the FSSP; it supports chronology, not a detailed sociology of recruitment.

The founding story still explained fidelity to Peter and the approved patrimony, but it could no longer describe every later member's personal passage. Formal formation and public institutional history therefore transmitted 1988 as inherited identity as well as living memory. That is a structural inference from the passage of time, not a measured account of candidate motives. A society born as a rescue and reconciliation structure was becoming a stable provider of formation and pastoral life in its own right.

\subsection{The Commission's end}

Francis decommissioned the Pontifical Commission \emph{Ecclesia Dei} in January 2019 and assigned its tasks to a section of the Congregation for the Doctrine of the Faith. The act said former \emph{Ecclesia Dei} institutes had achieved stability of number and life and that the remaining questions were predominantly doctrinal. For the FSSP this was both recognition of maturity and loss of the original supervisory home created at its birth.

The transfer did not suppress the society or its proper law. It changed the Roman office through which supervision occurred. Two years later \emph{Traditionis custodes} changed both the universal liturgical framework and the competent dicastery for the former \emph{Ecclesia Dei} societies.

\statusframe{Benedict XVI, \emph{Summorum Pontificum}, 7 July 2007; Francis, motu proprio decommissioning PCED, 17 January 2019.}{The general access regime broadened and then the original commission ended after the institutes had achieved institutional stability.}{Neither act is an FSSP erection or suppression. General liturgical law and the society's particular law must be distinguished at each date.}
